\section{Key-Value Cache Optimization}
\label{sec:kvcache}

The Key-Value cache represents one of the most critical bottlenecks in large language model inference, fundamentally constraining the deployment of long-context models in production environments. During autoregressive generation, the KV cache stores attention keys and values for all previously processed tokens, enabling efficient inference by avoiding redundant computation of past token representations. However, this optimization comes at a substantial memory cost. The cache size scales according to $2 \cdot n_{\text{layers}} \cdot n_{\text{heads}} \cdot d_{\text{head}} \cdot e \cdot b \cdot l$, where $e$ represents element size in bytes, $b$ denotes batch size, and $l$ indicates sequence length. For a LLaMA 3 70B model processing a 128K context window, this formulation yields approximately 40 GB of memory for a single request, dwarfing the model's parameter memory requirements.

The severity of this bottleneck extends beyond raw memory consumption. LLM inference operates in a fundamentally memory-bandwidth-bound regime rather than a compute-bound one. Generating each new token requires loading the entire KV cache from GPU memory, transforming what should be a computationally intensive operation into an I/O-dominated process. Empirical measurements of existing serving systems reveal utilization rates of only 20--38\% for allocated KV cache memory, with 60--80\% wasted through fragmentation and over-reservation to accommodate worst-case sequence lengths \citep{kwon2023efficient}. The inference efficiency plummets from approximately 70\% during training to merely 10\% during autoregressive generation, as the architecture shifts from compute-intensive parallel processing to memory-bandwidth-limited sequential token generation. This dramatic efficiency gap motivates a diverse ecosystem of optimization techniques spanning multiple dimensions: eviction policies that selectively discard tokens, quantization methods that reduce numerical precision, architectural modifications that fundamentally reduce cache requirements, and memory management systems that eliminate fragmentation waste.

\subsection{KV Cache Eviction Policies}

Eviction policies operate on the fundamental observation that not all cached tokens contribute equally to attention computation. By selectively retaining only the most influential tokens, these methods achieve substantial memory reductions while preserving model quality. The challenge lies in accurately identifying which tokens matter most, a determination that varies across attention heads, transformer layers, and generation contexts.

\subsubsection{H$_2$O: Heavy-Hitter Oracle}

\citet{zhang2024h2o} introduced H$_2$O (Heavy-Hitter Oracle), a pioneering eviction policy grounded in the empirical observation that attention values exhibit extreme concentration. Through systematic analysis across multiple model families, the authors demonstrated that a small subset of tokens---termed ``heavy hitters''---receives disproportionate attention weight, contributing the majority of effective value during attention computation. These heavy hitters emerge naturally during pre-training, correlating strongly with token co-occurrence patterns in natural language. Crucially, their removal results in catastrophic performance degradation, whereas selectively retaining them preserves model quality even at aggressive compression ratios.

H$_2$O formulates token selection as a dynamic submodular optimization problem, balancing retention of two categories: recent tokens to maintain position-aware context, and heavy-hitter tokens to preserve semantic importance. The method provides theoretical guarantees on the optimality of its selection strategy, ensuring that the chosen subset maximizes expected attention value under budget constraints. When evaluated on OPT-6.7B and OPT-30B models, H$_2$O achieves remarkable efficiency gains of up to 29$\times$ throughput improvement over DeepSpeed Zero-Inference and HuggingFace Accelerate, and up to 1.9$\times$ latency reduction compared to standard implementations. With only 20\% of tokens retained as heavy hitters, the method maintains model quality across diverse benchmarks spanning question answering, summarization, and language understanding tasks. The technique operates as a plug-and-play solution requiring no model fine-tuning, though it incurs computational overhead during the prefill phase to identify heavy hitters through attention profiling.

\subsubsection{ScissorHands and the Persistence of Importance}

Building on similar intuitions about token importance, \citet{liu2023scissorhands} proposed ScissorHands, which maintains the KV cache at a fixed budget through a different lens. The method rests on the ``persistence of importance'' hypothesis: pivotal tokens that substantially influenced attention at one generation step will continue to exert significant influence in future steps. This temporal stability of importance enables ScissorHands to make eviction decisions based on historical attention patterns rather than repeated profiling. The system stores pivotal tokens with higher retention probability, achieving up to 5$\times$ reduction in KV cache memory usage without requiring model fine-tuning or architectural modifications.

The effectiveness of ScissorHands extends further when combined with quantization techniques. Through a hybrid approach incorporating 4-bit quantization alongside selective token retention, the method achieves compression ratios up to 20$\times$ while maintaining competitive model quality. This combination demonstrates that eviction and quantization operate on complementary dimensions---eviction reduces the number of cached tokens, while quantization reduces the bits per token---enabling multiplicative rather than merely additive gains. The assumption of importance stability, however, introduces limitations. In scenarios where attention patterns shift dramatically during generation, such as multi-topic conversations or tasks requiring significant context switches, the persistence hypothesis may break down, resulting in suboptimal token selection.

\subsubsection{StreamingLLM and Attention Sinks}

\citet{xiao2024streamingllm} discovered a remarkable phenomenon that fundamentally reshapes our understanding of attention patterns in pre-trained language models. Their StreamingLLM framework revealed the existence of ``attention sinks''---initial tokens that receive consistently high attention scores despite lacking semantic importance. This phenomenon emerges from the softmax operation's constraint that attention weights must sum to one. When no token in the context provides a strong semantic match for a query, the model allocates excess attention weight to initial tokens, effectively using them as an attention ``sink'' or padding mechanism.

The practical implications of this discovery proved transformative for streaming applications. By retaining the KV states of merely four initial tokens alongside a sliding window of recent tokens, StreamingLLM enables LLMs to perform stable language modeling with effectively infinite sequence lengths. Models including LLaMA-2, MPT, Falcon, and Pythia successfully process sequences up to 4 million tokens without fine-tuning, achieving speedups of up to 22.2$\times$ compared to sliding window baselines that must recompute attention from scratch. The attention sink phenomenon emerges naturally during pre-training on sufficient data, becoming more pronounced in larger models and with extended training. Importantly, StreamingLLM operates as a training-free solution, requiring no position encoding modifications or interpolation techniques. The method has achieved widespread adoption, with integration into production systems including NVIDIA TensorRT-LLM and HuggingFace Transformers, demonstrating its practical utility for chatbots, document processing, and interactive systems requiring long-term context maintenance.

Subsequent theoretical work has illuminated the mechanisms underlying attention sink emergence. \citet{wu2025whenattentionsinkemerges} provided systematic analysis of when and why sinks form during training, demonstrating that they emerge after sufficient exposure to diverse linguistic patterns and become stronger with increased model capacity. The phenomenon reflects a learned strategy for handling the softmax constraint in the absence of semantically relevant tokens, suggesting fundamental implications for attention mechanism design.

\subsubsection{Adaptive Methods: FastGen, SnapKV, and PyramidKV}

Recognition that different attention heads and layers exhibit fundamentally different patterns motivated the development of adaptive eviction strategies. \citet{ge2024model} introduced FastGen, which employs targeted attention profiling during the prefill stage to discern the intrinsic structure of each attention module. Rather than applying a uniform eviction policy across all heads, FastGen adaptively selects from a portfolio of strategies: evicting long-range contexts for heads emphasizing local patterns, discarding non-special tokens for heads centered on special tokens like sentence boundaries, and employing standard full caching only for heads that broadly attend across all tokens. This adaptive approach achieves up to 55\% latency reduction at 16K generation length and 44.9\% pruned ratio on LLaMA 1-65B, with negligible generation quality loss compared to full caching. The method operates as a plug-and-play technique requiring no resource-intensive fine-tuning, profiling attention patterns once during prefill and applying the learned strategy throughout decoding.

\citet{li2024snapkv} extended the concept of head-specific selection with SnapKV, which selects a fixed number of tokens per attention head by retaining those with highest estimated contribution to the attention distribution. The method approximates token importance by accumulating attention weights received from an observation window of recent input tokens, smoothed with a moving average to reduce variance from stochastic generation. Building on this foundation, \citet{cai2024pyramidkv} introduced PyramidKV, which incorporates a critical insight from analysis of layer-wise attention patterns: deeper layers require less KV cache than shallower layers. The method assigns layer-specific retention budgets with larger allocations to lower layers, which focus on local pattern recognition and require broader context, and smaller allocations to upper layers, which perform higher-level semantic operations on more abstract representations. PyramidKV preserves performance metrics comparable to full caching using just 12\% of the standard KV cache (2048 tokens) on LongBench evaluations, and maintains acceptable quality with only 0.7\% of the full cache in extreme memory-constrained scenarios. The layer-aware budget allocation proves particularly effective for long-context tasks, where the pyramid structure naturally aligns with the information funneling that occurs across transformer depth.

\subsubsection{Q-HITTER and Quantization-Aware Selection}

A subtle but critical interaction between eviction and quantization motivated \citet{kang2024qhitter} to develop Q-HITTER, which reveals a fundamental limitation of combining H$_2$O-style selection with low-bit quantization. The problem manifests as a cascading error: H$_2$O identifies heavy hitters based on attention patterns computed with full-precision keys and values, but these selected tokens may possess activation distributions that are inherently difficult to quantize. High-magnitude outliers in heavy-hitter tokens cause quantization error to propagate through attention computation, corrupting the very patterns that motivated their selection. Furthermore, applying quantization before heavy-hitter identification distorts attention scores, causing selection to optimize for the wrong objective.

Q-HITTER addresses this coupling through dual-metric selection that jointly considers both attention importance and quantization friendliness. For each token, the method computes a composite score combining attention weight accumulation with a measure of how amenable the token's activation distribution is to quantization. Tokens with smooth, well-behaved distributions score higher on quantization friendliness than those with extreme outliers, even if their attention importance is comparable. This quantization-aware selection achieves near-optimal quality at 4-bit precision and below, where H$_2$O experiences significant degradation. Performance evaluations demonstrate impressive gains: up to 20$\times$ memory reduction with 33$\times$ throughput improvement over HuggingFace Accelerate, 7$\times$ over DeepSpeed, 4$\times$ over FlexGen, and 1.3$\times$ over H$_2$O itself, highlighting the importance of considering interactions between complementary optimization techniques.

\subsubsection{HashEvict and Locality-Sensitive Hashing}

Traditional attention-based eviction methods face a computational dilemma: identifying important tokens requires computing attention scores, but the goal of eviction is to avoid such computation in the first place. \citet{liu2024hashevict} proposed HashEvict to resolve this circularity through locality-sensitive hashing (LSH). Rather than computing full attention to determine token importance, the method employs LSH to rapidly approximate similarity between query tokens and cached key tokens. Binary hash codes derived from random projections enable fast Hamming distance computation, identifying dissimilar tokens for eviction before attention computation begins. This pre-attention decision making eliminates computational overhead: tokens evicted via LSH never enter the attention computation, yielding cumulative savings across layers.

Evaluated on LLaMA 3 across question answering, retrieval, and free-response tasks, HashEvict achieves 30--70\% KV cache compression with minimal performance degradation. The prefill phase demonstrates particularly impressive speedups of 1.5--2$\times$ compared to H$_2$O and ScissorHands, which must compute full attention before making eviction decisions. The method trades perfect similarity measurement for computational efficiency, relying on LSH's theoretical guarantees that close vectors in high-dimensional space map to similar hash codes with high probability. This approximation proves sufficient for eviction decisions, where the goal is identifying tokens that are clearly unimportant rather than perfectly ranking all tokens by importance.

\subsection{KV Cache Quantization}

Quantization reduces the precision of stored key and value tensors, directly decreasing both memory footprint and memory bandwidth requirements. The fundamental challenge lies in maintaining attention accuracy despite reduced numerical precision, as quantization errors propagate through the attention mechanism and accumulate across decoding steps. Successful quantization requires careful consideration of activation distributions, outlier patterns, and the sensitivity of attention computation to numerical noise.

\subsubsection{KIVI: Asymmetric 2-Bit Quantization}

\citet{liu2024kivi} developed KIVI, a tuning-free asymmetric 2-bit quantization method that achieves remarkable compression through a fundamental insight: keys and values exhibit different outlier patterns and thus benefit from different quantization strategies. Through systematic analysis of activation distributions across layers and attention heads, the authors discovered that key tensors exhibit high-magnitude outliers concentrated in specific channels, while value tensors display more uniform distributions with token-level rather than channel-level structure. This observation motivated an asymmetric quantization approach: keys receive per-channel quantization to handle channel-wise outliers, while values receive per-token quantization aligned with streaming inference patterns.

The asymmetric strategy enables KIVI to operate at aggressive 2-bit precision while maintaining model quality comparable to full precision inference. Evaluations on LLaMA, Falcon, and Mistral model families demonstrate 2.6$\times$ reduction in peak memory usage (including model weights), enabling 4$\times$ larger batch sizes for the same memory budget. Throughput improvements range from 2.35$\times$ to 3.47$\times$ on real LLM workloads, achieved through carefully optimized CUDA kernels that efficiently quantize and dequantize during inference. The per-token quantization strategy for values proves particularly beneficial for streaming inference, where new tokens append directly to the existing cache without requiring re-quantization of previous tokens. This streaming compatibility eliminates a major source of computational overhead in naive quantization approaches. KIVI operates as a plug-and-play solution requiring zero fine-tuning, making it immediately applicable to pre-trained models and enabling rapid deployment in production serving systems.

\subsubsection{KVQuant and Extreme Context Lengths}

Targeting extreme context lengths up to 10 million tokens, \citet{hooper2024kvquant} introduced KVQuant, a comprehensive quantization framework incorporating four novel techniques that synergistically address the challenges of sub-4-bit quantization. First, per-channel key quantization adjusts quantization dimensions to match the inherent structure of key activation distributions, grouping elements along the channel dimension where outliers naturally concentrate. Second, pre-RoPE key quantization applies quantization before rotary positional embeddings, mitigating the distributional shift that RoPE introduces and enabling more accurate quantized representations. Third, non-uniform KV cache quantization derives per-layer sensitivity-weighted datatypes, recognizing that different layers tolerate quantization differently and allocating precision accordingly. Fourth, per-vector dense-and-sparse quantization explicitly isolates numerical outliers on a per-vector basis, preventing extreme values from skewing quantization ranges and degrading the representation of typical values.

This multi-pronged approach achieves less than 0.1 perplexity degradation with 3-bit quantization on WikiText-2 and C4 benchmarks, enabling previously impossible inference scenarios: LLaMA-7B processing a 1 million token context on a single A100-80GB GPU, and 10 million token contexts on an 8-GPU system with 8$\times$ KV cache compression. Custom CUDA kernels optimized for low-precision matrix-vector multiplication deliver 1.7$\times$ speedups at 4-bit precision compared to FP16 baseline operations. The non-uniform quantization scheme proves particularly important at extreme context lengths, where the cumulative effect of layer-wise quantization errors would otherwise cause catastrophic quality degradation. By carefully tuning per-layer precision based on sensitivity analysis, KVQuant maintains the delicate balance between memory efficiency and representational capacity required for million-token inference. The method has achieved integration into vLLM, one of the most widely deployed LLM serving frameworks, demonstrating its production readiness and practical utility.

\subsubsection{GEAR: Hybrid Quantization with Low-Rank Correction}

\citet{kang2024gear} proposed GEAR, recognizing that pure quantization approaches face fundamental limitations: at very low bit-widths, quantization alone cannot maintain accuracy across the diversity of activation patterns in KV cache. GEAR addresses this through a three-component hybrid approach. The majority of KV cache entries, which exhibit similar magnitude scales, receive ultra-low-precision 4-bit quantization. Quantization residuals---the errors introduced by quantization---are captured through low-rank matrix approximation via singular value decomposition, exploiting the observation that residual errors often lie in a low-dimensional subspace. Individual outlier errors that escape both quantization and low-rank approximation receive explicit sparse matrix correction, handling extreme values through a separate sparse representation that maintains high precision for outliers while applying aggressive compression to typical values.

This hybrid strategy achieves near-lossless 4-bit compression with up to 2.38$\times$ throughput improvement and 2.29$\times$ peak memory reduction. The method explicitly addresses a challenge unique to autoregressive generation: errors compound at each decoding step, as the corrupted KV cache influences the generation of subsequent tokens, which in turn adds to the cache and propagates errors forward. By maintaining a low-rank correction term that captures systematic quantization biases, GEAR prevents the accumulation of correlated errors that would otherwise cause quality drift during long sequence generation. The sparse correction component handles the statistically rare but semantically critical outliers that often carry outsized importance for model behavior, ensuring that compression does not inadvertently discard the very information that determines generation quality.

\subsubsection{KITTY and Information-Theoretic Quantization}

Recent work by \citet{yuan2024kitty} approaches 2-bit quantization through the lens of information theory, proposing KITTY (K-means Information Theory-guided Token-wise quantization for accurate 2-bit KV cache). The method employs clustering algorithms to group tokens with similar activation patterns, then assigns quantization codebooks optimized for each cluster. Information-theoretic analysis guides the clustering process, ensuring that quantization preserves the mutual information between full-precision and quantized representations. This token-wise clustering approach adapts to the diversity of activation patterns across different contexts, allocating representational capacity where it provides maximal information preservation. Preliminary evaluations suggest improved accuracy compared to uniform 2-bit quantization at comparable computational overhead, though the method requires additional profiling to determine optimal cluster assignments.

\subsection{Architectural Modifications}

While eviction and quantization optimize within the constraints of standard transformer architecture, architectural modifications fundamentally reduce KV cache requirements by altering how attention operates. These changes typically require training or uptraining but deliver permanent efficiency gains that compound with other optimization techniques.

\subsubsection{Multi-Query Attention}

\citet{shazeer2019fast} introduced Multi-Query Attention (MQA) as a direct response to the memory bandwidth bottleneck in autoregressive decoding. Standard Multi-Head Attention maintains separate key and value projections for each attention head, requiring the model to load $H$ copies of key-value tensors during inference, where $H$ denotes the number of attention heads. Shazeer's analysis identified this repeated loading of similar information as the primary performance limiter on modern accelerators. MQA modifies the architecture to share a single set of key and value projections across all query heads, while maintaining separate query projections to preserve representational diversity. This sharing dramatically reduces KV cache size proportional to the number of attention heads, typically yielding 8--16$\times$ reduction for models with 8--16 attention heads.

The memory bandwidth savings translate directly to inference speedup, as each token generation requires loading the shared key-value state once rather than $H$ times. Quality evaluations revealed minimal accuracy degradation from this architectural change, particularly when models train with MQA from initialization. The technique enables substantially larger batch sizes within fixed memory budgets, improving throughput for serving scenarios. Early adopters including LLaMA-2 and Falcon demonstrated MQA's effectiveness in production models, though subsequent work would reveal that the single shared key-value head sometimes introduces a quality-efficiency trade-off that later methods would seek to better balance.

\subsubsection{Grouped-Query Attention}

\citet{ainslie2023gqa} proposed Grouped-Query Attention (GQA) as a generalization interpolating between standard Multi-Head Attention and Multi-Query Attention. GQA divides the $H$ query heads into $G$ groups, with each group sharing a single key-value head. This formulation recovers MHA when $G = H$ (each group contains one head) and MQA when $G = 1$ (all heads form a single group), while intermediate values of $G$ explore the trade-off space between memory efficiency and representational capacity. The authors demonstrated that GQA captures much of MQA's efficiency benefit while mitigating quality degradation, particularly important for larger models where the capacity reduction from single key-value head becomes more pronounced.

A critical contribution of the GQA work is the uptraining recipe: converting existing MHA checkpoints to GQA requires only 5\% of original pre-training compute through a fine-tuning procedure that leverages the knowledge already captured in MHA weights. This low conversion cost enabled rapid adoption, as organizations could retrofit existing models rather than training from scratch. The quality-efficiency trade-off proves favorable across model scales, with GQA achieving near-MHA quality at speeds comparable to MQA. Following publication in May 2023, GQA achieved remarkably swift adoption in production models. Meta incorporated GQA in LLaMA 2 (July 2023) and LLaMA 3 (2024), while Mistral AI employed it in Mistral-7B (September 2023). By 2024, GQA had become the de facto standard for new model releases, demonstrating the practical importance of balanced KV cache reduction.

For concrete illustration, consider LLaMA 3 8B: with standard MHA, processing a long context requires approximately 280 GB of KV cache, triggering out-of-memory errors on most available GPUs. With GQA configured to use 4 groups, the same model requires approximately 70 GB of KV cache---a 75\% reduction that enables inference on consumer hardware while maintaining model quality competitive with the MHA baseline. This dramatic improvement in memory requirements fundamentally expands the deployment scenarios for large language models, making long-context inference practical on single-GPU systems.

\subsubsection{Cross-Layer KV Sharing}

\citet{brandon2024reducing} introduced Cross-Layer Attention (CLA), extending the sharing concept beyond heads within a layer to sharing across adjacent layers. The method recognizes that transformer layers form a compositional hierarchy, with each layer performing increasingly abstract transformations of the input representation. CLA computes key-value projections for only a subset of layers, allowing subsequent layers to reuse KV activations from previous layers rather than computing fresh projections. At both 1B and 3B model scales, applying CLA with a sharing factor of 2 (computing KV for every other layer) yields 2$\times$ KV cache reduction while incurring less than 1\% perplexity degradation for MQA models with typical head dimensions of 64 and 128. Crucially, CLA operates orthogonally to both MQA and GQA, enabling combination with either technique for multiplicative rather than additive gains.

\citet{yang2024layercondensed} systematically investigated cross-layer sharing through their Layer-Condensed KV Cache (LCKV) framework, which uniformly divides transformer layers into groups of adjacent layers and pairs queries from all layers in each group with KVs from the group's bottom layer. When reducing KV cache size by 2$\times$, most LCKV configurations achieve higher throughput than standard transformers while maintaining competitive performance. The systematic study by \citet{tu2024systematic} further demonstrated that cross-layer sharing enables 2$\times$ throughput improvements across diverse model architectures and sizes, with careful analysis of which sharing patterns prove most effective for different model families and training regimes. The authors found that sharing works better in middle-to-deep portions of models, where adjacent layers produce increasingly similar representations, compared to early layers where representations change rapidly or final layers where output-specific transformations occur.

\subsubsection{MiniCache and Depth-Dimension Compression}

\citet{jia2024minicache} developed MiniCache, exploiting the high similarity between KV states in adjacent layers at middle-to-deep portions of transformer models. Rather than discarding less important layers' KV states entirely, MiniCache compresses across the depth dimension through magnitude-direction decomposition. The method disentangles each KV state into a magnitude component (the norm of the activation vector) and a direction component (the unit-normalized vector). For adjacent layers with highly similar directions, MiniCache interpolates the directional components while preserving the distinct magnitudes, achieving compression through the observation that directions change slowly across layers while magnitudes may vary. Token retention mechanisms preserve highly distinct state pairs unmerged, minimizing information loss for the minority of tokens where adjacent-layer representations diverge significantly.

On ShareGPT dialogues, LLaMA-2-7B with 4-bit MiniCache achieves 5.02$\times$ compression, approximately 5$\times$ throughput enhancement, and 41\% memory footprint reduction compared to FP16 full cache baseline. The depth-dimension compression proves complementary to sequence-dimension techniques like token eviction, as they operate on orthogonal axes of the KV cache tensor. A system could potentially combine PyramidKV's layer-wise token budgeting with MiniCache's cross-layer compression, evicting unimportant tokens within each layer while merging similar representations across layers, though such combined approaches remain an open research direction.

\subsection{Memory Management Innovations}

Beyond reducing cache size, innovative memory management addresses fragmentation waste and enables efficient sharing of KV cache across requests. These systems-level optimizations operate orthogonally to compression techniques, providing multiplicative efficiency gains when combined.

\subsubsection{PagedAttention}

\citet{kwon2023efficient} introduced PagedAttention, drawing direct inspiration from virtual memory and paging in operating systems. Traditional attention implementations require contiguous memory for keys and values, but sequences have unpredictable lengths that necessitate either pre-allocating maximum-length buffers (wasting memory for short sequences) or dynamic allocation (causing fragmentation as sequences of different lengths finish at different times). PagedAttention resolves this through block-based allocation: the KV cache partitions into fixed-size blocks, each containing keys and values for a fixed number of tokens (typically 16). Contiguous logical blocks in a sequence map to potentially non-contiguous physical blocks through per-sequence block tables that translate logical to physical addresses, analogous to page tables in OS virtual memory.

This design achieves near-optimal memory utilization with under 4\% waste compared to 60--80\% in existing systems, as only the final block of each sequence contains unused slots. The block table mechanism enables sophisticated memory management strategies. Multiple sequences can share physical blocks containing identical content, critical for scenarios like beam search where hypotheses share common prefixes, or batched inference where multiple users submit requests with common prompt templates. Blocks deallocate immediately upon sequence completion and enter a free pool for allocation to new requests, eliminating the fragmentation that plagues traditional contiguous allocation schemes.

The vLLM serving system built on PagedAttention achieves 2--4$\times$ throughput improvement compared to FasterTransformer and Orca while maintaining equivalent latency. The gains compound at larger batch sizes and longer sequences, precisely the operating regime where KV cache memory becomes most constrained. PagedAttention has achieved widespread adoption as the foundation for modern LLM serving systems, with deployment in major cloud providers and integration into open-source frameworks. The success of PagedAttention demonstrates that careful systems engineering, applying decades-old OS concepts to new domains, can deliver order-of-magnitude improvements even in heavily optimized systems.

\subsubsection{vAttention: OS-Level Demand Paging}

\citet{liu2024vattention} proposed vAttention as an alternative approach to dynamic memory management, questioning whether application-level paging is necessary when operating systems already provide sophisticated demand paging mechanisms. PagedAttention requires custom attention kernels aware of block tables and capable of gathering from non-contiguous physical blocks. These custom kernels introduce maintenance burden and potential performance overhead. vAttention instead leverages OS and GPU support for virtual memory: KV cache occupies contiguous virtual address space, but the operating system lazily allocates physical pages only when accessed, automatically evicts pages under memory pressure, and transparently handles the mapping between virtual and physical memory.

This design preserves the abstraction of contiguous memory for attention kernels, enabling use of highly optimized standard attention implementations without modification. The OS handles paging decisions based on access patterns, potentially capturing workload-specific locality better than application-level heuristics. Evaluated on LLaMA models, vAttention achieves 1.97$\times$ faster token generation compared to vLLM's PagedAttention, and 3.92$\times$ faster prompt processing compared to PagedAttention's FlashAttention variant. The method maintains near-zero memory waste comparable to PagedAttention while eliminating custom kernel complexity.

However, vAttention faces deployment constraints: not all GPU architectures provide unified virtual memory with demand paging support, and OS page sizes may not align optimally with attention computation granularity. These hardware dependencies limit vAttention's applicability compared to PagedAttention's broader compatibility. The tension between these two approaches reflects a broader question in systems design: whether domain-specific application-level management or general-purpose OS-level mechanisms better optimize for specialized workloads. Current evidence suggests both have roles, with PagedAttention dominating deployment due to wider hardware support, while vAttention demonstrates intriguing possibilities for future hardware-software co-design.

\subsubsection{Prompt Caching}

\citet{gao2024prompt} addressed a different dimension of memory management: reuse of computation across requests sharing common prompt segments. Many production workloads exhibit significant prompt overlap: system messages repeat across all requests from an application, retrieval-augmented generation systems incorporate the same documents into multiple queries, and conversation systems maintain shared context across turns. Standard inference recomputes attention for these shared segments in every request, wasting both computation and memory bandwidth.

Prompt Cache introduces a schema-based approach where users explicitly mark reusable prompt modules through special tokens. During initial processing, the system computes and caches attention states (KV representations) for marked modules, associating them with unique identifiers. Subsequent requests referencing the same modules fetch cached states rather than recomputing, assembling a complete KV cache from a mix of cached and newly-computed components. This modular attention reuse delivers dramatic latency improvements: 8--60$\times$ reduction in time-to-first-token depending on whether execution occurs on GPU (8$\times$) or CPU (60$\times$), with the larger CPU gains reflecting the higher cost of attention computation on less specialized hardware.

The technique proves particularly impactful for chatbot systems with shared system prompts, retrieval-augmented generation with repeated document contexts, multi-turn conversations reusing history, and batch inference with common instructions. Major API providers including OpenAI and Anthropic have deployed prompt caching in production, offering reduced per-token pricing for cached content and thereby incentivizing users to structure prompts for reusability. The explicit schema requirement represents both a strength and limitation: users must consciously design for caching through proper prompt structure, but this explicitness ensures cache correctness and predictable behavior compared to automatic schemes that might incorrectly share states across semantically different contexts.

\subsubsection{SqueezeAttention and Layer-Wise Budget Allocation}

\citet{li2024squeezeattention} introduced SqueezeAttention, optimizing KV cache allocation across two dimensions simultaneously: the sequence dimension (which tokens to keep) and the layer dimension (how much cache per layer). Traditional methods optimize only sequence dimension, applying uniform per-layer budgets. SqueezeAttention recognizes that transformer layers exhibit heterogeneous sensitivity to KV cache size, with some layers maintaining quality under aggressive compression while others degrade rapidly. The method measures per-layer importance through cosine similarity between input perturbations and output changes, identifying layers that merely pass through information (low importance) versus layers that substantially transform representations (high importance).

Based on importance scores, SqueezeAttention reallocates cache budgets from less sensitive to more sensitive layers. A layer that tolerates 70\% token eviction with minimal degradation contributes its excess budget to a sensitive layer that requires 90\% retention for quality preservation. This redistribution achieves 30--70\% memory reduction with up to 2.2$\times$ throughput improvement, outperforming uniform budget allocation that must choose between aggressive compression everywhere (degrading quality on sensitive layers) or conservative compression (wasting capacity on tolerant layers). The two-dimensional optimization proves particularly effective for long-context inference, where the cumulative effect of per-layer budget decisions determines overall system efficiency.

\subsection{Integrated Systems and Future Directions}

The diversity of KV cache optimization techniques raises important questions about integration and composition. \citet{kang2024entropy} proposed entropy-guided dynamic budget allocation, using attention entropy as a signal for per-layer cache requirements. Layers with high attention entropy---indicating broad, diffuse attention across many tokens---receive larger budgets, while layers with concentrated attention patterns tolerate smaller caches. On Qwen3-4B, this approach achieves 4.18\% memory reduction while preserving ROUGE scores, and on Mistral-0.1v-7B delivers 46.6\% decoding time reduction. The entropy metric provides a theoretically grounded alternative to empirical importance measurement, though computing attention entropy introduces overhead that must be amortized across sufficient generation steps.

Research into combining complementary techniques remains nascent. Q-HITTER demonstrated the importance of considering quantization-eviction interactions, but comprehensive frameworks that jointly optimize eviction, quantization, architectural choices, and memory management remain underdeveloped. \citet{paul2024pitfalls} provided critical analysis of KV cache compression limitations, identifying scenarios where compression fails: out-of-distribution inputs that violate importance assumptions, reasoning-heavy tasks sensitive to information loss, and long-horizon dependencies that aggressive compression disrupts. The work by \citet{wang2024holdonto} specifically examined reasoning task sensitivity, finding that compression affects multi-step reasoning disproportionately compared to simple question answering, suggesting task-specific compression strategies may be necessary.

Long-context generalization methods like LM-Infinite \citep{han2024lminfinite} and InfLLM \citep{zhang2024infllm} explore alternative architectural approaches. LM-Infinite employs lambda-shaped attention masks and distance limits to achieve training-free generalization from 2K to 32K tokens with 7.5$\times$ memory reduction and 2.7$\times$ decoding speedup. InfLLM introduces memory-based long context handling, storing distant context in a separate memory structure with token-relevant lookup, enabling inference on sequences up to 1 million tokens without expanding the attention window. These methods suggest that fundamental rethinking of attention mechanisms, beyond optimization within standard transformer architecture, may unlock further efficiency gains.

The YOCO architecture \citep{sun2024yoco} exemplifies radical architectural redesign. Rather than the standard self-attention decoder, YOCO employs a decoder-decoder architecture with a self-attention encoder followed by cross-attention decoder. The self-attention encoder processes tokens and produces a global KV cache, which all cross-attention decoder layers share, reducing memory scaling from $O(N \times L)$ for standard transformers to $O(N)$, where $N$ denotes sequence length and $L$ denotes layer count. This architectural change delivers approximately 4$\times$ memory reduction, though it requires training from scratch rather than uptraining from existing checkpoints.

FlashAttention \citep{dao2022flashattention} and its successors optimize attention computation through IO-aware algorithms that minimize memory transfers between GPU memory hierarchy levels. While primarily targeting training efficiency, FlashAttention reduces inference memory footprint by up to 20$\times$ and achieves 3$\times$ speedup for sequences of 128--2K tokens, with continued efficiency gains up to 64K tokens. The IO-aware perspective proves complementary to KV cache compression: FlashAttention optimizes how we compute attention given a fixed cache size, while compression techniques reduce the cache size itself.

Industry deployment reveals the practical importance of KV cache optimization. NVIDIA's technical documentation on CPU-GPU memory sharing for KV cache offload demonstrates hardware-software co-design approaches, leveraging NVLink and PCIe bandwidth to store portions of KV cache in CPU memory when GPU memory saturates, trading bandwidth for capacity. The LMCache system \citep{liu2024lmcache} addresses enterprise-scale serving across multiple nodes, incorporating fault tolerance, cost optimization, and cache coherence protocols for distributed KV cache. These systems-level concerns, often overlooked in research focused on single-model optimization, prove critical for production deployment at scale.

The rapid pace of innovation in KV cache optimization throughout 2023--2024 reflects both the critical importance of this bottleneck and the substantial headroom for further improvement. Multiple research directions remain open: unified frameworks combining complementary techniques, hardware-aware optimization tailored to specific accelerator architectures, dynamic adaptation adjusting compression strategies based on input characteristics, extreme context lengths pushing toward 100M+ tokens, and extension to multi-modal models where KV cache must accommodate visual tokens alongside text. The compression ratios of 2--20$\times$ demonstrated by current techniques, often with minimal quality degradation, suggest that KV cache management will remain a productive research area as model sizes and context requirements continue their upward trajectory. As deployment scenarios diversify---from edge devices to data center clusters, from chatbots to complex reasoning systems---the portfolio of available optimization techniques grows increasingly valuable, enabling practitioners to select and combine methods appropriate for their specific constraints and objectives.
